\documentclass[aps,prfluids,preprint,superscriptaddress,nofootinbib]{revtex4-2}

\usepackage{amsmath}
\usepackage{amssymb}
\usepackage{graphicx}
\usepackage{booktabs}
\usepackage{array}
\usepackage{tabularx}
\usepackage{fvextra}
\fvset{breaklines=true,fontsize=\footnotesize}
\graphicspath{{../}}
\newcommand{\code}[1]{\texttt{#1}}

\begin{document}

\title{Two clocks, not one: measured pressure--temperature decoupling, its operator origin, and the structural blind spot of single-diffusivity closure}

\author{Abdelrahman Saifelden}
\email[Correspondence: ]{Abdulrahman.Saifelden22011411@aiet.edu.eg}
\affiliation{Alexandria Higher Institute of Engineering and Technology (AIET), Alexandria, Egypt}

\date{\today}

\begin{abstract}
Pressure and temperature in a fluid are driven by the same source yet evolve under different operators: temperature obeys a local, parabolic advection--diffusion equation, while pressure in the incompressible limit obeys a global elliptic constraint --- the finite-Mach hyperbolic acoustic adjustment degenerating into it as $M\to 0$. This paper assembles the direct observational and computational case that this operator split is measurable as \emph{two clocks} running in one flow, and that single-eddy-diffusivity (``K-theory'') closure is structurally blind to it. \textbf{Temporally}: in NEON eddy-covariance data (WREF, 30-min, winter and summer months) temperature is a clean diurnal (24 h) signal while barometric pressure is synoptic-plus-tidal, and the bulk daily $P$--$T$ correlation is regime-dependent but never gas-law-locked ($r\approx +0.07$ in winter, $-0.49$ in the summer thermal-low regime, against a constant-density lock slope of $+0.33$ kPa/K); in ASOS 1-minute data at eight stations spanning $24.6^{\circ}$--$47.4^{\circ}$N the 12-h pressure tide amplitude falls monotonically $1.10\to 0.34$ hPa toward the pole (daily-block-bootstrap 95\% CIs; trend correlation $0.98$ with the Haurwitz/Dai--Wang migrating-tide climatology $1.16\cos^{3}\varphi$, station totals $0.9$--$1.4\times$ the migrating-only formula as non-migrating components require) while the 24-h temperature amplitude stays local --- pressure carries a planetary-scale resonance where temperature carries the sun overhead; band-resolved coherence completes the split: $0.8$--$1.0$ at the solar lines (the shared astronomical clock, above a phase-randomized surrogate 95\% level at all eight stations) collapsing to the surrogate noise floor in the off-line sub-diurnal band. \textbf{Spatially}: in a $Ra=10^{6}$ Rayleigh--B\'enard DNS the two fields share one integral scale (ratio $1.06\times$) yet the pressure Taylor microscale is $2.40\times$ the buoyancy's and the high-wavenumber pressure spectrum sits $\sim 10^{4}\times$ below buoyancy --- the $1/k^{2}$ footprint of the inverse Laplacian; sweeping $Ra=10^{6}\to 10^{8}$ \emph{widens} the microscale ratio to $4.5\times$ while the shared integral scale stays order one, so the split strengthens with turbulence intensity; correspondingly, a localized boundary bump projects 79\% of its elliptic pressure response beyond one bump-width while the matched parabolic response stays confined. The \textbf{closure-relevant consequence} is measured directly: across Monin--Obukhov stability classes the momentum-to-heat transport-efficiency ratio swings $4.3\times$ ($0.77\to 3.28$) in the winter month and $6.0\times$ in the summer month, structurally outside any fixed-turbulent-Prandtl single-diffusivity closure. Three \textbf{controlled demonstrations} isolate the mechanism: (i) a compressible pressure field relaxes to the elliptic Poisson field with $t_{90}\approx 2.3\,L/c$, all sound speeds collapsing onto one curve in $ct/L$; (ii) in a fully nonlinear compressible box the dilatational (fast-clock) energy fraction scales as $M^{2}$ and the acoustic-averaged pressure converges to the elliptic field (correlation $\approx 1$, residual $\sim M^{2}$); (iii) the projection method is exhibited as the two-clocks structure made literal --- one Leray solve reduces the local drift's divergence from $2.7\times 10^{-2}$ to machine zero ($2\times 10^{-15}$), whereas phase-randomized surrogates with the \emph{exactly} matched power spectrum leave it at $3.7\times 10^{-2}$ [$3.5$--$3.8\times 10^{-2}$, 95\% over $n=100$ surrogates] --- no better than applying no correction at all ($2.7\times 10^{-2}$) --- with pointwise correlation to the true correction of median $|r|=0.07$ (95th percentile $0.20$; the seed ensemble replaces a single-realization anecdote) --- a structural ceiling of spectrum-matching stochastic closures. Finally, the same decomposition applied to \textbf{real winds} (NCEP/NCAR reanalysis, spherical-harmonic Helmholtz split) yields $KE_{div}/KE_{rot}=4.0\%,\,0.7\%,\,1.0\%$ at 850/500/250 hPa with the minimum at the classical level of non-divergence --- reproduced in the independent NCEP--DOE Reanalysis 2 ($4.5\%,\,0.8\%,\,1.1\%$, same ordering, same 500-hPa minimum; $\le 15\%$ relative differences), so the ratio is not an assimilation-system artifact --- and daily averaging scrubs the fast clock exactly as acoustic averaging does in the box. Scope is explicit: the demonstrations are two-dimensional and periodic; no 3-D regularity claim and no forecast-skill claim is made; reanalysis is a model-assimilated confirmation at $l\le 35$. A companion paper derives, audits, and benchmarks the closure repair; this paper establishes the phenomenon it repairs.
\end{abstract}

\maketitle

\section*{1. Introduction}

\subsection*{1.1 One source, two responses}

Turn on the sun over a fluid layer and two of its most basic state variables respond in ways that could hardly be more different. Temperature responds \emph{locally}: it is advected, it diffuses, it remembers the surface that heated it, and its anomalies are torn into filaments by whatever strain the flow supplies. Pressure responds \emph{globally}: in the incompressible limit it is not a prognostic field at all but the multiplier of the divergence-free constraint, resolved instantaneously over the whole domain and its boundaries; at finite compressibility that instantaneous adjustment is replaced by acoustic waves running at a speed far faster than the flow. The two fields are coupled --- through buoyancy, through the equation of state, through the momentum balance --- and driven by the same heating; yet they carry different histories, different geometries, and different characteristic times. We call this the \emph{two clocks} structure: a slow, local, memory-bearing clock (temperature and other transported scalars) and a fast, global, constraint-like clock (pressure).

The distinction is elementary at the level of the equations. It stops being elementary the moment one builds a turbulence closure, a subgrid parameterization, or a statistical forecast model, because the workhorse of all three --- a single scalar eddy diffusivity acting on resolved gradients (Boussinesq 1877; Smagorinsky 1963), optionally with a fixed turbulent Prandtl number to relate momentum and heat --- assumes \emph{one} clock. Everything transported is assumed stirred by the same eddies at the same rate with no memory and no action at a distance. The premise of this paper is that the two-clocks structure is not a metaphor but a measurable property of real data and simple flows, with a quantifiable failure mode for one-clock closures.

\subsection*{1.2 What this paper does}

We proceed from measurement to mechanism to consequence:

\begin{enumerate}
\item \textbf{Measure the temporal split in the field} (\S3). NEON eddy-covariance data and multi-station 1-minute ASOS barometry/thermometry show temperature locked to the local diurnal cycle while pressure is dominated by synoptic variability plus the semidiurnal solar tide --- a \emph{global} resonance whose amplitude follows latitude, not local heating. Bulk daily $P$--$T$ coupling at a fixed site is regime-dependent and never gas-law-locked (winter $r\approx +0.07$, summer $r\approx -0.49$).
\item \textbf{Measure the spatial split on one grid} (\S4). In a Rayleigh--B\'enard DNS, pressure and buoyancy share the same integral (roll) scale but differ $2.40\times$ in Taylor microscale and $\sim 10^{4}\times$ in high-$k$ spectral power --- the direct footprint of the inverse Laplacian's $1/k^{2}$; the microscale ratio \emph{widens} to $4.5\times$ as $Ra$ is swept $10^{6}\to 10^{8}$, and a Green's-function experiment makes the same point with a single localized source.
\item \textbf{Measure the closure-relevant consequence} (\S5). Binned by Monin--Obukhov stability, the ratio of momentum to heat transport efficiency swings $4.3\times$ where a fixed-$Pr_{t}$ closure requires it constant.
\item \textbf{Demonstrate the mechanism in controlled solvers} (\S6). A linear acoustics solver shows the hyperbolic$\to$elliptic crossover with $t_{90}\approx 2.3\,L/c$; a nonlinear compressible solver shows both clocks coexisting with the fast clock's energy $\sim M^{2}$; a Boussinesq solver shows that the projection method \emph{is} the two-clocks synthesis, and that a spectrum-matched stochastic surrogate cannot reproduce the elliptic correction (median $|r|=0.07$, p95 $0.20$, $n=100$; divergence unrepaired) --- the structural ceiling of SPDE-style closures.
\item \textbf{Confirm on real winds} (\S7). A spherical-harmonic Helmholtz split of NCEP/NCAR reanalysis puts $\lesssim 4\%$ of kinetic energy in the divergent (fast) component, with the vertical structure pinned to the classical level of non-divergence and time-averaging acting as the same low-pass filter as acoustic averaging in the box.
\end{enumerate}

Section 8 states scope and kill conditions for each claim. A companion manuscript (Saifelden 2026, hereafter \textbf{P1}) takes the diagnosis to prescription: it audits eddy-diffusivity closure operator-by-operator against the Mori--Zwanzig generalized Langevin equation and benchmarks a projected, fluctuation--dissipation-consistent repair in 2-D and 3-D. This paper is the evidence layer: every number above is regenerable from committed code (Appendix A), and none of it depends on P1's closure machinery.

\subsection*{1.3 Relation to existing understanding}

None of the ingredients is exotic. That low-Mach-number flow relaxes to an elliptically-constrained state is classical singular-limit theory (Klainerman \& Majda 1981); atmospheric tides are a mature subject (Chapman \& Lindzen 1970; Dai \& Wang 1999); the rotational dominance of large-scale winds and the level of non-divergence are textbook synoptic dynamics (Holton 2004); projection methods are standard numerics (Chorin 1968; Leray 1934); and stochastic parameterization is an active program (Hasselmann 1976; Berner et al. 2017). The contribution here is the \emph{joint, quantitative} case built for the closure question: the same operator split measured in field data, in a DNS, in controlled compressible and incompressible solvers, and in reanalysis, each with a falsifiable number attached, and each mapped onto the specific term of the closure problem it breaks. In particular the spectrum-matching ceiling of \S6.3 --- energy-consistent noise reproducing the fast clock's \emph{spectrum} while leaving its \emph{geometry} unreconstructed (ensemble-median correlation $0.07$, the divergence constraint unrepaired) --- is, to our knowledge, not stated quantitatively elsewhere, and it is the piece that connects the phenomenology to closure structure.

\medskip\hrule\medskip

\section*{2. Two operators, two clocks}

\subsection*{2.1 The equations}

Let $u$ be velocity, $p$ pressure, $\theta$ a transported scalar (temperature, buoyancy, humidity), $\nu,\kappa$ molecular coefficients. The scalar obeys a \textbf{parabolic} transport equation,
\[
\partial_{t}\theta + (u\cdot\nabla)\theta = \kappa\nabla^{2}\theta + S_{\theta},
\tag{1}
\]
whose Green's function is (in the diffusive limit) a spreading Gaussian of width $\sqrt{4\kappa t}$: \emph{local}, causal, memory-bearing. In an incompressible flow, pressure obeys the \textbf{elliptic} Poisson constraint obtained from the divergence of the momentum equation,
\[
\nabla^{2}p = -\rho\,\nabla\!\cdot\!\left[(u\cdot\nabla)u\right] + \nabla\!\cdot\! f,
\tag{2}
\]
whose Green's function is $\sim\ln r$ in 2-D ($\sim -1/r$ in 3-D): \emph{global}, instantaneous, boundary-aware. In Fourier space $\hat{p}=-\widehat{f}/k^{2}$ --- division by $k^{2}$ is a global low-pass filter, and this single algebraic fact drives every spatial measurement in \S4. At finite compressibility pressure is instead \textbf{hyperbolic},
\[
\frac{1}{c^{2}}\,\partial_{tt}p - \nabla^{2}p = s(x,t),
\tag{3}
\]
with sound speed $c$; as $M=U/c\to 0$, (3) degenerates into (2) --- the waves become infinitely fast and vanish into the constraint (Klainerman \& Majda 1981). The two descriptions are the same physics at different compressibility, which is exactly what \S6.1--6.2 demonstrate numerically.

\subsection*{2.2 The clocks}

Three timescales matter for a domain of size $L$, flow speed $U$, eddy scale $\ell$:
\[
\tau_{p}\sim L/c \;\;(\text{acoustic/elliptic adjustment}), \qquad
\tau_{adv}\sim \ell/U \;\;(\text{advective turnover}), \qquad
\tau_{\kappa}\sim \ell^{2}/\kappa \;\;(\text{diffusion}).
\tag{4}
\]
For air at laboratory-to-mesoscale $L$, $\tau_{p}/\tau_{adv}=M\,(L/\ell)^{-1}\ll 1$; for water ($c\approx 1500$ m/s) smaller still. Pressure equilibrates before the flow moves; scalars integrate their history as the flow moves. A closure that models both with one local diffusivity therefore commits to a specific, testable set of errors:

\begin{itemize}
\item \textbf{E1 (temporal).} At a fixed site, low-frequency pressure and temperature should decouple: temperature tracks the local source spectrum; pressure tracks global/planetary modes. A one-clock picture in which ``$P$ tracks $T$'' through the gas law at fixed density predicts tight daily coupling --- so measuring the coupling tests the picture. (\S3)
\item \textbf{E2 (spatial).} On one grid, pressure must lack small-scale power relative to a scalar by the $1/k^{4}$ spectral factor implied by $\hat{p}\propto \hat{f}/k^{2}$, even where both share the energy-containing scale. (\S4)
\item \textbf{E3 (closure).} Momentum (shaped by the global pressure field) and heat (shaped by local buoyant plumes) cannot maintain a constant transport-efficiency ratio as the balance of shear and buoyancy changes; a fixed-$Pr_{t}$ closure says they must. (\S5)
\item \textbf{E4 (crossover).} Making $c$ finite must turn the constraint into waves whose only role, in the slow limit, is to set the adjustment clock: relaxation time $\propto L/c$, final field $\to$ the elliptic solution, fast-clock energy $\to 0$ as $M^{2}$. (\S6.1--6.2)
\item \textbf{E5 (structure, not spectrum).} The fast clock's correction to the slow dynamics is a global elliptic solve, not noise: a stochastic surrogate with the correct spectrum but random phases must fail to reproduce it pointwise and must fail its constraint function. (\S6.3)
\item \textbf{E6 (real flow).} In the real atmosphere the slow clock (rotational, balanced) must dominate the kinetic energy, the fast clock (divergent) must be weak and physically structured, and time-averaging must act as a low-pass filter on it. (\S7)
\end{itemize}

Each expectation is falsifiable, and \S8.2 lists what observation would have killed --- or would still kill --- each one.

\medskip\hrule\medskip

\section*{3. Temporal decoupling measured in the field}

\subsection*{3.1 NEON eddy covariance: a winter month at one tower, cross-checked in summer}

We use the NEON bundled eddy-covariance product (DP4.00200.001) at WREF (Wind River Experimental Forest, WA; 45.82$^{\circ}$N, 121.95$^{\circ}$W; old-growth conifer canopy, 53 m), January 2020 at 30-min resolution: 1488 intervals, of which 1097 carry valid pressure and 1180 valid temperature. Three facts (Fig.~1):

\begin{itemize}
\item \textbf{The source--response chain is causal and ordered.} Incoming shortwave leads air temperature by $\sim$2 h and the sensible heat flux by $\sim$3 h; daytime flux is upward (buoyant). The temperature clock is driven, locally, by the sun.
\item \textbf{The two spectra are qualitatively different.} Temperature variance is overwhelmingly diurnal (24 h). Pressure variance is dominated by multi-day synoptic weather plus a distinct \textbf{semidiurnal (12 h) atmospheric tide} --- twin daily maxima visible in the raw composite.
\item \textbf{Bulk coupling is near zero.} Pearson $r(P,T)\approx +0.07$ ($R^{2}\approx 0.005$) at the daily scale: nowhere near the steep constant-density line a fixed-parcel ideal-gas lock would predict. The mismatch is absorbed by $\sim$1\% density changes --- the gas law holds \emph{locally}; what fails is the one-clock reading ``bulk $P$ tracks bulk $T$''.
\end{itemize}

\begin{center}
\includegraphics[width=0.86\linewidth]{figures/01_diurnal_composites.png}

{\small \textbf{Figure 1.} NEON WREF, Jan 2020: diurnal composites of shortwave, temperature, sensible heat flux, and pressure. Temperature is a single clean 24-h signal; pressure shows the twin-peak semidiurnal tide riding on synoptic variability.}
\end{center}

\textbf{Seasonal cross-check (summer month).} The same pipeline applied to July 2020 at the same tower (1036/1344 valid $P$/$T$ intervals) reproduces the fingerprint at the opposite solar extreme: temperature's semidiurnal-to-diurnal power ratio drops to $1.0\%$ (the summer diurnal line is even cleaner; winter $7.3\%$), pressure keeps a genuine tide line at $11.9$ h ($8.5\%$ of its diurnal power) with its dominant variability still synoptic ($\sim$168 h), and the bulk daily correlation moves to $r=-0.49$ --- the summer thermal-low regime, in which hot afternoons \emph{lower} surface pressure --- with observed slope $-0.024$ kPa/K against the constant-density lock's $+0.334$ kPa/K. The bulk $P$--$T$ coupling is set by synoptic regime and season; the operator fingerprint is not (artifact \code{01b\_neon\_seasonal.json}).

\subsection*{3.2 ASOS 1-minute: the latitude fingerprint of a global mode}

A single tower cannot distinguish ``pressure is noisy'' from ``pressure listens to the whole atmosphere''. Eight ASOS stations spanning $23^{\circ}$ of latitude, coastal and interior (1-minute data, Q1 2020; Iowa Environmental Mesonet archive; daily moving-block bootstrap, $n=500$, for the 95\% CIs) can:

\begin{center}\small
\begin{tabular}{lccccc}
\toprule
station & latitude & $S_{1,T}$ (24 h) & $S_{2,P}$ (12 h) [95\% CI] & $S_{2,P}/1.16\cos^{3}\varphi$ & daily $P$--$T$ corr [95\% CI] \\
\midrule
EYW & 24.6$^{\circ}$N & 1.51 $^{\circ}$C & \textbf{1.103 [1.013, 1.191] hPa} & 1.26 & $-0.16$ $[-0.30, -0.03]$ \\
MIA & 25.8$^{\circ}$N & 2.73 $^{\circ}$C & 1.077 [0.979, 1.169] hPa & 1.27 & $-0.05$ $[-0.18, +0.08]$ \\
MSY & 30.0$^{\circ}$N & 2.75 $^{\circ}$C & 1.018 [0.878, 1.169] hPa & 1.35 & $-0.40$ $[-0.54, -0.24]$ \\
DFW & 32.9$^{\circ}$N & 3.28 $^{\circ}$C & 0.957 [0.794, 1.141] hPa & 1.39 & $-0.48$ $[-0.59, -0.34]$ \\
PHX & 33.4$^{\circ}$N & 5.06 $^{\circ}$C & 0.924 [0.861, 1.008] hPa & 1.37 & $-0.20$ $[-0.29, -0.12]$ \\
DSM & 41.5$^{\circ}$N & 2.73 $^{\circ}$C & 0.628 [0.398, 0.965] hPa & 1.29 & $-0.56$ $[-0.67, -0.41]$ \\
CAR & 46.9$^{\circ}$N & 3.03 $^{\circ}$C & 0.439 [0.181, 0.879] hPa & 1.19 & $-0.42$ $[-0.52, -0.30]$ \\
SEA & 47.4$^{\circ}$N & 1.79 $^{\circ}$C & \textbf{0.337 [0.262, 0.529] hPa} & 0.94 & $-0.06$ $[-0.21, +0.09]$ \\
\bottomrule
\end{tabular}
\end{center}

The 24-h temperature amplitude is set by local insolation and does not order with latitude in this winter-season sample (interior PHX largest, marine SEA/EYW smallest). The 12-h pressure tide \textbf{falls monotonically from equator to pole across all eight stations}, the classical signature of the global solar semidiurnal tide --- a planetary normal-mode response, largest where the forcing projects onto the low-latitude waveguide (Chapman \& Lindzen 1970). This is deliberately a \textbf{protocol validation against a known planetary signal}, not a discovery claim: the trend correlation with the migrating-tide climatology $A_{H}=1.16\cos^{3}\varphi$ hPa (Haurwitz 1956) is $0.98$, and station totals sit $0.9$--$1.4\times$ $A_{H}$, the expected excess once non-migrating components are included (Dai \& Wang 1999). Pressure ``knows about'' the geometry of the whole atmosphere; temperature knows about the sun overhead (Figs.~2--3). Station-level daily $P$--$T$ correlations range $-0.05$ to $-0.56$ with 95\% CIs excluding any tight positive lock --- variable in sign and magnitude, set by synoptic regime rather than by any gas-law coupling. \textbf{Band-resolved coherence} sharpens the bulk correlation (the blunt instrument) into the actual statement: magnitude-squared coherence between hourly $P$ and $T$ is $0.82$--$0.98$ \emph{at} the solar lines ($S_{1}$/$S_{2}$, where both fields are phase-locked to the same astronomical forcing; above the phase-randomized-surrogate 95\% level at all eight stations) and collapses to the surrogate noise floor ($0.09$--$0.30$ vs floor $0.15$--$0.20$) in the off-line sub-diurnal band (9--11 h and 13--20 h): away from the shared solar clock, nothing locks the two fields.

\begin{center}
\includegraphics[width=0.95\linewidth]{figures/11b_tide_vs_latitude_ci.png}

{\small \textbf{Figure 2.} Left: semidiurnal pressure-tide amplitude vs latitude at eight stations with bootstrap 95\% CIs, against the Haurwitz $1.16\cos^{3}\varphi$ migrating-tide climatology: monotonic equator-to-pole decay --- a global resonance, not local heating. Right: band coherence --- high at the shared solar lines, at the surrogate floor off-line.}
\end{center}

\begin{center}
\includegraphics[width=0.8\linewidth]{figures/12_asos_coupling.png}

{\small \textbf{Figure 3.} Daily-scale pressure--temperature coupling at the three ASOS stations: weak and regime-dependent, nowhere near a constant-density gas-law line.}
\end{center}

\textbf{Scope.} Two months at the seasonal extremes (NEON) and one season (ASOS) at nine sites in one country: these are \emph{fingerprint} measurements --- the diurnal/semidiurnal split and its latitude ordering are the tested claims --- not a global climatology. The tide amplitudes agree with the published climatological maps (Dai \& Wang 1999) where they overlap.

\medskip\hrule\medskip

\section*{4. Spatial decoupling measured on one grid}

\subsection*{4.1 Rayleigh--B\'enard DNS: same forcing, two geometries}

Field sensors cannot separate spatial structure; a DNS can. We use the 2-D Rayleigh--B\'enard configuration of The Well (Ohana et al. 2024): $Ra=10^{6}$, $Pr=1$, $512\times 128$, periodic in $x$, no-slip walls, with velocity, buoyancy, and pressure co-resolved on one grid. The decoupling is \emph{not} about dominant sizes --- both fields share the convection-roll wavelength (integral lengths 0.195 vs 0.207, ratio $1.06\times$). It lives at small scales (Figs.~4--5):

\begin{center}
\begin{tabular}{lccc}
\toprule
metric & buoyancy $\theta$ & pressure $p$ & ratio $p/\theta$ \\
\midrule
integral length (large scale) & 0.195 & 0.207 & $1.06\times$ \\
Taylor microscale (small scale) & 0.253 & 0.607 & $\mathbf{2.40\times}$ \\
high-$k$ spectral power ($k\approx 10$) & --- & --- & $\theta$ larger by $\sim\mathbf{10^{4}\times}$ \\
\bottomrule
\end{tabular}
\end{center}

Both numbers are the direct footprint of \S2: $\hat{p}\propto\hat{f}/k^{2}$ suppresses spectral power by $\sim 1/k^{4}$, so pressure is smooth and blobby (elliptic) while buoyancy is advected into thin filaments with intense small-scale gradients (parabolic). Nothing was tuned: the same coupled simulation produces both fields.

\textbf{Robustness across Ra.} The single-snapshot $2.40\times$ is not a one-$Ra$ accident. Sweeping The Well's $Pr=1$ test files at $Ra=10^{6},10^{7},10^{8}$ (5 trajectories $\times$ 6 statistically-steady snapshots each) gives Taylor-microscale ratios of $2.58$ [$1.98$--$3.59$], $2.94$ [$2.54$--$3.65$], and $4.45$ [$3.27$--$5.88$] (median [5--95\%]; the headline snapshot sits inside the $Ra=10^{6}$ interval), while the integral-length ratio stays order one ($1.0$--$1.8$) and the high-$k$ spectral gap widens by another $\sim 70\times$. As $Ra$ grows and buoyancy plumes sharpen toward finer scales, the inverse Laplacian keeps holding pressure at the large scales: \emph{the elliptic--parabolic split widens with turbulence intensity rather than closing} (artifact \code{15b\_rb\_ra\_sweep.json}).

\begin{center}
\includegraphics[width=0.86\linewidth]{figures/13_rb_fields.png}

{\small \textbf{Figure 4.} RB DNS snapshot: buoyancy (filamentary, sharp) vs pressure (smooth, domain-spanning) on the same grid at the same instant.}
\end{center}

\begin{center}
\includegraphics[width=0.8\linewidth]{figures/15_rb_spectra.png}

{\small \textbf{Figure 5.} Spatial spectra: shared energy-containing scale; $\sim 10^{4}\times$ high-$k$ gap --- the $1/k^{4}$ elliptic suppression.}
\end{center}

\subsection*{4.2 A single localized source: Green's functions in action}

To isolate the operator contrast from turbulence statistics, we add a single localized bump to the bed of a developed cavity flow (frozen field), compute the continuity source it induces, $q=-\nabla\!\cdot\!(\text{bump}\cdot u)$, and apply that \emph{same} source to both operators. The elliptic response $\Delta p = \nabla^{-2}q$ spreads over the whole cavity --- \textbf{79\%} of its magnitude lies beyond one bump-width of the source --- while the heat-kernel (parabolic) response at matched time remains a confined Gaussian blob. Boundary geometry edits the \emph{global} pressure field instantly, and only the \emph{local} temperature field: any closure that models the pressure-borne flux with a local diffusivity is using the wrong Green's function.

\begin{center}
\includegraphics[width=0.8\linewidth]{figures/39_elliptic_pressure.png}

{\small \textbf{Figure 6.} Identical localized source, two operators: global $\ln r$ elliptic response (79\% beyond one bump-width) vs confined Gaussian parabolic response.}
\end{center}

\medskip\hrule\medskip

\section*{5. The closure-relevant consequence, measured}

If momentum is shaped by the global elliptic clock and heat by the local parabolic clock, then the \emph{efficiencies} with which turbulence transports the two cannot stay in a fixed ratio as the shear/buoyancy balance changes. First-order K-theory with a constant turbulent Prandtl number ($K_{M}\approx Pr_{t}^{-1}K_{H}$) says they must.

Using the same NEON tower (602 usable 30-min periods), we bin the transport-efficiency correlations $r_{uw}$ (momentum) and $|r_{wT}|$ (heat) by Monin--Obukhov stability $\zeta=(z-d)/L$ (Monin \& Obukhov 1954; Foken 2008):

\begin{center}
\begin{tabular}{lcccc}
\toprule
stability class & $n$ & median $r_{uw}$ & median $|r_{wT}|$ & momentum/heat ratio \\
\midrule
strongly unstable & 95 & 0.182 & 0.209 & 0.87 \\
unstable & 73 & 0.324 & 0.150 & 2.17 \\
near-neutral & 121 & 0.312 & 0.095 & \textbf{3.28} \\
stable & 153 & 0.225 & 0.200 & 1.13 \\
strongly stable & 160 & 0.156 & 0.203 & \textbf{0.77} \\
\bottomrule
\end{tabular}
\end{center}

Momentum efficiency peaks near-neutral (shear-driven, pressure-mediated); heat efficiency peaks under strong convection (buoyant plumes) --- they move \emph{oppositely}, and their ratio swings $\mathbf{4.3\times}$ ($0.77\to 3.28$). The summer month repeats the swing with more convective sampling: $0.43\to 2.60$ across the same $\zeta$ classes ($6.0\times$; 1030 usable periods), so the blind spot is not a winter artifact. The stability dependence itself is classical --- turbulent-Prandtl-number variation across $\zeta$ is embedded in the Businger--Dyer flux--profile relations (Businger et al. 1971) and reviewed at length by Li (2019) --- so the row-level numbers validate the protocol against known surface-layer physics; the claim we \emph{use} is only the structural one that follows: A single fixed-$Pr_{t}$ diffusivity is structurally unable to reproduce a $4.3\times$ swing in this ratio; stability-dependent similarity functions absorb it empirically, which is precisely the point --- the physics forces the closure to grow a second, stability-indexed degree of freedom because one clock is not enough. This is E3, and it is the field-data anchor for the operator audit in P1.

\medskip\hrule\medskip

\section*{6. Controlled demonstrations of the mechanism}

\subsection*{6.1 The hyperbolic $\to$ elliptic crossover (linear)}

A minimal 2-D linear-acoustics solver makes the elliptic-vs-hyperbolic statement concrete. An initial pressure bump radiates as a ring at speed $c$ and decays by geometric spreading (the wave regime). Under a fixed forcing, the compressible pressure relaxes to the \emph{same} field as the elliptic Poisson solve --- final relative error $\sim 2\times 10^{-3}$ for every $c$ tested --- and the approach time scales as
\[
t_{90}\approx 2.3\,L/c ,
\tag{5}
\]
with error histories for all sound speeds collapsing onto a single curve in the rescaled time $ct/L$ (Fig.~7). $c$'s only role is to set the clock; as $c\to\infty$ the adjustment becomes instantaneous and (3) degenerates into (2). This is why water ($c\approx 1500$ m/s) is ``more fiercely elliptic'' than air ($c\approx 340$ m/s).

\begin{center}
\includegraphics[width=0.8\linewidth]{figures/17_crossover.png}

{\small \textbf{Figure 7.} Crossover: relaxation of the compressible pressure to the elliptic field; all $c$ collapse in $ct/L$; $t_{90}\approx 2.3\,L/c$.}
\end{center}

\subsection*{6.2 Both clocks in one nonlinear flow}

A fully nonlinear 2-D isothermal compressible Navier--Stokes solver (pseudo-spectral, periodic, $p=c^{2}\rho$ so $c$ is an explicit knob) is initialized from a Taylor--Green vortex at $Re\approx 600$ and Helmholtz-split into solenoidal (vortical, slow) and dilatational (acoustic, fast) parts:

\begin{center}
\begin{tabular}{ccc}
\toprule
Mach $M$ & $\langle KE_{dil}\rangle/\langle KE_{sol}\rangle$ & avg-pressure residual vs elliptic \\
\midrule
0.05 & $1.5\times 10^{-4}$ & $8.1\times 10^{-4}$ \\
0.10 & $5.9\times 10^{-4}$ & $3.2\times 10^{-3}$ \\
0.20 & $2.3\times 10^{-3}$ & $1.4\times 10^{-2}$ \\
0.40 & $9.3\times 10^{-3}$ & $1.0\times 10^{-1}$ \\
\bottomrule
\end{tabular}
\end{center}

Two clocks coexist in one flow: the vortical energy decays on the slow advective time while the acoustic energy oscillates on the fast period (Fig.~8), and the fast clock's energy fraction scales as $\sim M^{2}$, vanishing in the incompressible limit (Fig.~9). The \emph{instantaneous} pressure is dominated by the standing acoustic wave; averaging over the fast clock (the quasi-steady second half of each run) scrubs $\tau_{p}$ and leaves a field matching the elliptic Poisson pressure at correlation $\approx 1$, with the residual itself shrinking as $M^{2}$. Averaging over the fast clock is a \emph{low-pass filter that recovers the elliptic limit} --- an operation \S7 will repeat on the real atmosphere.

\begin{center}
\includegraphics[width=0.82\linewidth]{figures/20_ns_two_clocks.png}

{\small \textbf{Figure 8.} Nonlinear compressible box: slow vortical decay with fast acoustic oscillation riding on it --- both clocks, one flow.}
\end{center}

\begin{center}
\includegraphics[width=0.8\linewidth]{figures/21_ns_mach_sweep.png}

{\small \textbf{Figure 9.} Mach sweep: dilatational energy fraction and acoustic-averaged pressure residual both $\sim M^{2}$.}
\end{center}

\subsection*{6.3 The synthesis: projection is the two-clocks structure --- and spectra cannot fake it}

The standard projection (fractional-step) method advances a Boussinesq flow by a slow drift followed by an elliptic correction:
\[
u^{*}=u^{n}+\Delta t\left[-(u\cdot\nabla)u+\nu\nabla^{2}u+(b-\langle b\rangle)\hat{y}\right],\qquad
u^{n+1}=\mathbb{P}u^{*},\qquad
b^{n+1}=b^{n}+\Delta t\left[-(u\cdot\nabla)b+\kappa\nabla^{2}b\right],
\tag{6}
\]
with $\mathbb{P}=I-\nabla(\nabla^{2})^{-1}\nabla\cdot$ the Leray projector (Leray 1934; Chorin 1968). Read structurally: the drift is the generator of the local Markov/heat semigroup --- the \emph{perfect} local, memoryless forecaster, exactly what single-K closure and statistical predictors represent --- and the projection is one global Poisson solve per step, the fast clock acting as a constraint. Three results (Figs.~10--11):

\begin{itemize}
\item \textbf{The convection cell reappears from first principles}: plumes rise under the slow clock, the elliptic clock pulls in the compensating inflow (entrainment), and the shear between them rolls into vortices.
\item \textbf{The local drift is divergent; one elliptic solve fixes it instantly.} The intermediate $u^{*}$ violates mass conservation at RMS $\nabla\!\cdot\!u^{*}\approx 2.7\times 10^{-2}$; a single projection drops it to $\approx 2\times 10^{-15}$ --- machine zero --- through the global potential $\phi$ with $\nabla^{2}\phi=\nabla\!\cdot\!u^{*}$.
\item \textbf{Energy yes, structure no.} Replacing the projection correction by a phase-randomized field with the \emph{identical} power spectrum (Parseval-exact) --- the structural move of spectrum-matched stochastic/SPDE closures (Hasselmann 1976; Berner et al. 2017) --- yields pointwise correlation with the true correction of median $|r|=0.07$ (95th percentile $0.20$ over a 100-surrogate ensemble; the correction's energy sits in a few large-scale modes, so single-draw correlations scatter --- we report the ensemble, not a favourable draw) and leaves RMS $\nabla\!\cdot\!u = 3.7\times 10^{-2}$ [$3.5$--$3.8\times 10^{-2}$, 95\%] instead of zero --- no better than applying no correction at all ($2.7\times 10^{-2}$). Matching a spectrum does not reconstruct a boundary-aware, globally-coupled elliptic operator.
\end{itemize}

\begin{center}
\includegraphics[width=0.82\linewidth]{figures/23_projection_step.png}

{\small \textbf{Figure 10.} One time step, dissected: divergent local drift $u^{*}$; the global elliptic potential; the projected, solenoidal update (RMS $\nabla\!\cdot\!u$: $2.7\times 10^{-2}\to 2\times 10^{-15}$).}
\end{center}

\begin{center}
\includegraphics[width=0.82\linewidth]{figures/24_spde_limit.png}

{\small \textbf{Figure 11.} The SPDE ceiling: a Parseval-exact phase-randomized surrogate of the projection correction --- correct spectrum, correlation median $|r|=0.07$ (p95 $0.20$, $n=100$), constraint unrepaired ($\nabla\!\cdot\!u\approx 3.7\times 10^{-2}$, 95\% interval $3.5$--$3.8\times 10^{-2}$).}
\end{center}

This is E5, and it is the sharpest structural statement in the paper: the fast clock is not ``unresolved variance'' to be re-injected with the right energy --- it is a global operator with a definite geometric job. Stochastic closures that match spectra can repair energetics; they cannot repair geometry. (P1 builds the repair that can: a \emph{projected}, dissipation-tied stochastic force --- the projection is what rescues the noise.)

\medskip\hrule\medskip

\section*{7. The two clocks in the real atmospheric wind}

Finally, the same decomposition on real winds. NCEP/NCAR Reanalysis 1 horizontal winds (2.5$^{\circ}$; Kalnay et al. 1996), fetched live over OPeNDAP, are Helmholtz-decomposed on the sphere via the spherical-harmonic inversions $\nabla^{2}\psi=\zeta$, $\nabla^{2}\chi=\delta$ into a rotational (non-divergent, balanced --- slow) and a divergent (fast) wind.

\textbf{Honesty note.} Horizontal divergence is not a model artifact: by 3-D continuity it is the genuine footprint of vertical motion. The diagnostic is therefore the energy split, not the existence of divergence.

\begin{center}
\begin{tabular}{lcc}
\toprule
level & R1: $KE_{div}/KE_{rot}$ (daily mean) & R2 cross-check \\
\midrule
850 hPa & 4.0\% & 4.5\% \\
500 hPa & \textbf{0.7\%} & \textbf{0.8\%} \\
250 hPa & 1.0\% & 1.1\% \\
\bottomrule
\end{tabular}
\end{center}

\begin{itemize}
\item \textbf{The weather is rotational; the fast clock is a weak, structured residual.} The rotational wind \emph{is} the synoptic highs/lows and jets; the divergent wind is $\sim$100$\times$ weaker in energy and concentrated in convergence/overturning zones --- tropical convection, storm inflow --- the real-atmosphere analogue of the dilatational component (Fig.~12).
\item \textbf{Divergent $\ll$ rotational at every scale.} Splitting kinetic energy by spherical wavenumber $l$, the divergent branch sits 1--2 decades below the rotational branch for all resolved $l$; integrated, $KE_{div}/KE_{rot}\approx 0.04$ at 850 hPa --- the real-data echo of $KE_{dil}/KE_{sol}\sim M^{2}\ll 1$ (Fig.~13).
\item \textbf{The vertical structure is physics, not tuning.} The divergent fraction is smallest near 500 hPa --- the classical \emph{level of non-divergence} (Holton 2004) --- and larger in the boundary layer and jet-outflow levels (Fig.~14). \textbf{Cross-checked in a second assimilation system:} recomputing the same window with the same Helmholtz split from NCEP--DOE Reanalysis 2 (different model physics and assimilation corrections; Kanamitsu et al. 2002) gives $4.5\%/0.8\%/1.1\%$ --- same ordering, same 500-hPa minimum, $\le 15\%$ relative differences (\code{run\_\allowbreak{}reanalysis\_\allowbreak{}r2.py}, \code{figures/\allowbreak{}27b\_\allowbreak{}reanalysis\_\allowbreak{}crosscheck.json}); the ratio is a property of the analyzed atmosphere, not of one reanalysis.
\item \textbf{Time-averaging is the low-pass filter.} The 6-hourly ratio at 500 hPa (1.1\%) exceeds the daily-mean ratio (0.7\%): averaging scrubs the fast divergent clock --- the same operation that recovered the elliptic pressure from the acoustic field in \S6.2.
\end{itemize}

\begin{center}
\includegraphics[width=0.86\linewidth]{figures/25_reanalysis_two_clocks.png}

{\small \textbf{Figure 12.} NCEP/NCAR winds, one analysis day: rotational wind (synoptic systems) vs divergent wind (weak, confined to overturning zones).}
\end{center}

\begin{center}
\includegraphics[width=0.8\linewidth]{figures/26_reanalysis_ke_spectrum.png}

{\small \textbf{Figure 13.} Kinetic-energy spectra by spherical wavenumber: the divergent branch 1--2 decades below the rotational at every $l\le 35$.}
\end{center}

\begin{center}
\includegraphics[width=0.8\linewidth]{figures/27_reanalysis_levels.png}

{\small \textbf{Figure 14.} Vertical structure of $KE_{div}/KE_{rot}$: minimum at the level of non-divergence; 6-hourly $>$ daily --- time-averaging as low-pass.}
\end{center}

\textbf{Scope.} Reanalysis is a model-assimilated product, not raw observation, and 2.5$^{\circ}$ resolves only $l\le 35$. This section is a real-data \emph{confirmation} of the energy split and its structure --- not a turbulence-closure claim and not a gravity-wave climatology.

\medskip\hrule\medskip

\section*{8. Discussion}

\subsection*{8.1 What has been established}

Six independent measurements, one structure. At a fixed site pressure and temperature run on different clocks with regime-dependent, never gas-law-locked bulk daily coupling (E1); on one DNS grid pressure lacks exactly the small scales the inverse Laplacian says it must lack (E2); the momentum/heat transport-efficiency ratio swings $4.3\times$ in winter and $6.0\times$ in summer where one-clock closure requires a constant (E3); finite compressibility turns the constraint into waves whose only slow-limit role is the clock $t_{90}\approx 2.3L/c$, with fast-clock energy $\sim M^{2}$ (E4); the fast clock's action on the slow dynamics is a global elliptic correction that spectrum-matched noise cannot imitate --- correlation $0.01$, constraint unrepaired (E5); and the real atmospheric wind carries $\lesssim 4\%$ of its energy in the fast clock, structured exactly where the physics puts it (E6).

\subsection*{8.2 Kill conditions (what would have falsified, or still would)}

\begin{center}
\begin{tabular}{p{0.08\linewidth}p{0.55\linewidth}p{0.28\linewidth}}
\toprule
claim & observation that kills it & status \\
\midrule
E1 & tight positive daily $P$--$T$ lock at fixed sites; tide amplitude ordered by local heating rather than latitude & not observed ($r\approx+0.07$ winter, $-0.49$ summer; tide $1.08\to0.63$ hPa poleward) \\
E2 & pressure microscale $\approx$ scalar microscale; no high-$k$ suppression & not observed ($2.40\times$; $10^{4}\times$) \\
E3 & transport-efficiency ratio constant ($\pm$tens of \%) across stability & not observed ($4.3\times$ swing) \\
E4 & relaxation time not $\propto L/c$; compressible field $\not\to$ Poisson field & not observed (collapse in $ct/L$) \\
E5 & spectrum-matched surrogate reproducing the projection correction pointwise & not observed (ensemble median $|r|=0.07$, p95 $0.20$; div unrepaired) \\
E6 & $KE_{div}\sim KE_{rot}$, or structureless divergent wind & not observed ($\le 4\%$, structured) \\
\bottomrule
\end{tabular}
\end{center}

Each claim remains open to stronger tests: more sites and seasons (E1), higher-$Ra$ and 3-D DNS (E2), other towers and surface types (E3), 3-D compressible sweeps (E4), wall-bounded domains where projection and filtering no longer commute (E5; see P1 \S5d), and higher-resolution reanalyses or raw soundings (E6).

\subsection*{8.3 What this means for closures --- and what it does not}

The measurements above do not say K-theory is useless; they say it is \emph{one-clock} physics, calibrated where the slow clock dominates and structurally silent about the fast one. Concretely: (i) any single-diffusivity model inherits E3's contradiction as soon as its calibration regime changes; (ii) any stochastic augmentation that matches spectra but not the elliptic geometry inherits E5's ceiling; (iii) any scheme that time-averages its inputs has already low-passed the fast clock and should not expect to recover fast-clock phenomena downstream (E4/E6). The constructive direction --- collapse of the Mori--Zwanzig memory kernel, restoration of the projected fluctuation--dissipation force, the resulting benchmark closures in 2-D and 3-D --- is the companion paper's subject (P1; cf. Chorin, Hald \& Kupferman 2000; Parish \& Duraisamy 2017). This paper deliberately claims none of it.

Equally deliberately, nothing here is a regularity statement (no Beale--Kato--Majda claim; the demonstrations are 2-D and periodic), nothing is an operational forecast-skill claim, and the reanalysis section stands on an assimilated product at coarse resolution. The claims are exactly the six numbers of \S8.1, at the stated scopes.

\section*{9. Conclusion}

Pressure and temperature share a source and a fluid but not a clock: one is a global elliptic/acoustic adjustment, the other a local parabolic transport. That split is measurable in a forest, at three airports, in a convection DNS, in two compressible solvers, in a projection step, and in the planetary wind field --- six falsifiable expectations, six confirmations, with quantitative kill conditions on record. Its practical bite is structural: one-diffusivity closures and spectrum-matched stochastic augmentations are blind to the fast clock's geometry, measurably so ($4.3\times$ efficiency-ratio swing; ensemble-median correlation $0.07$ of spectrum-matched surrogates with the true elliptic correction, divergence left unrepaired). The phenomenon established here is the load-bearing premise of the closure audit and repair developed in the companion paper.

\medskip\hrule\medskip

\section*{Appendix A --- Reproducibility}

All figures and numbers regenerate from the public repository (single commands, fixed seeds, CPU-only for every result in this paper):

\begin{center}
\begin{tabular}{p{0.30\linewidth}p{0.44\linewidth}p{0.16\linewidth}}
\toprule
result & command & figures \\
\midrule
\S3.1 NEON two clocks & \code{python general\_two\_clocks/run\_analysis.py} & 1 (and lead--lag, spectra) \\
\S3.1 seasonal cross-check & \code{python general\_two\_clocks/run\_neon\_compare.py} & --- (artifact) \\
\S3.2 ASOS latitude sweep + CIs & \code{python general\_two\_clocks/asos/fetch\_iem.py \&\& python general\_two\_clocks/run\_asos\_ci.py} & 15--20 \\
\S4.1 RB DNS split & \code{python general\_two\_clocks/run\_rb.py} & 4--5 \\
\S4.1 Ra sweep ($10^{6}$--$10^{8}$) & \code{python general\_two\_clocks/run\_rb\_ra\_sweep.py} & --- (artifact) \\
\S4.2 Green's-function demo & \code{python general\_two\_clocks/run\_elliptic\_demo.py} & 6 \\
\S5 stability bins & \code{python general\_two\_clocks/run\_stability.py} & --- (table) \\
\S6.1 linear crossover & \code{python general\_two\_clocks/run\_compressible.py} & 7 \\
\S6.2 nonlinear Mach sweep & \code{python general\_two\_clocks/run\_compressible\_ns.py} & 8--9 \\
\S6.3 projection/SPDE & \code{python general\_two\_clocks/run\_boussinesq.py} & 10--11 \\
\S7 reanalysis split & \code{python general\_two\_clocks/run\_reanalysis.py} & 12--14 \\
\bottomrule
\end{tabular}
\end{center}

Data endpoints: NEON DP4.00200.001 (data.neonscience.org, public); ASOS 1-minute via the Iowa Environmental Mesonet archive (mesonet.agron.iastate.edu, public); The Well RB dataset (polymathic-ai, public); NCEP/NCAR Reanalysis 1 via NOAA PSL OPeNDAP (psl.noaa.gov, public, no credentials). The repository's test suite (656 passing tests at draft time) covers the analysis modules; each report in \code{general\_two\_clocks/} documents per-part method details and robustness checks.

\section*{Data availability}

All analysis code, committed figure data, and the manuscript source are in the public repository at https://github.com/abdh0saifelden2-debug/K. All input datasets are public as listed in Appendix A.

\section*{References}

Berner, J., et al. (2017). Stochastic parameterization: toward a new view of weather and climate models. \emph{Bulletin of the American Meteorological Society} 98, 565--588. doi:10.1175/BAMS-D-15-00268.1

Boussinesq, J. (1877). Essai sur la th\'eorie des eaux courantes. \emph{M\'emoires pr\'esent\'es par divers savants \`a l'Acad\'emie des Sciences} 23, 1--680.

Businger, J. A., Wyngaard, J. C., Izumi, Y., \& Bradley, E. F. (1971). Flux-profile relationships in the atmospheric surface layer. \emph{Journal of the Atmospheric Sciences} 28, 181--189.

Chapman, S., \& Lindzen, R. S. (1970). \emph{Atmospheric Tides: Thermal and Gravitational}. D. Reidel, Dordrecht.

Chorin, A. J. (1968). Numerical solution of the Navier--Stokes equations. \emph{Mathematics of Computation} 22, 745--762. doi:10.1090/S0025-5718-1968-0242392-2

Chorin, A. J., Hald, O. H., \& Kupferman, R. (2000). Optimal prediction and the Mori--Zwanzig representation of irreversible processes. \emph{Proceedings of the National Academy of Sciences} 97, 2968--2973. doi:10.1073/pnas.97.7.2968

Dai, A., \& Wang, J. (1999). Diurnal and semidiurnal tides in global surface pressure fields. \emph{Journal of the Atmospheric Sciences} 56, 3874--3902.

Foken, T. (2008). \emph{Micrometeorology}. Springer, Berlin.

Hasselmann, K. (1976). Stochastic climate models. Part I: Theory. \emph{Tellus} 28, 473--485. doi:10.1111/j.2153-3490.1976.tb00696.x

Haurwitz, B. (1956). The geographical distribution of the semidiurnal solar pressure oscillation. \emph{New York University Meteorological Papers} 2(5), 36 pp.

Holton, J. R. (2004). \emph{An Introduction to Dynamic Meteorology}, 4th ed. Elsevier Academic Press.

Kalnay, E., et al. (1996). The NCEP/NCAR 40-year reanalysis project. \emph{Bulletin of the American Meteorological Society} 77, 437--471.

Kanamitsu, M., Ebisuzaki, W., Woollen, J., Yang, S.-K., Hnilo, J. J., Fiorino, M., \& Potter, G. L. (2002). NCEP--DOE AMIP-II Reanalysis (R-2). \emph{Bulletin of the American Meteorological Society} 83, 1631--1643. doi:10.1175/BAMS-83-11-1631

Klainerman, S., \& Majda, A. (1981). Singular limits of quasilinear hyperbolic systems with large parameters and the incompressible limit of compressible fluids. \emph{Communications on Pure and Applied Mathematics} 34, 481--524. doi:10.1002/cpa.3160340405

Leray, J. (1934). Sur le mouvement d'un liquide visqueux emplissant l'espace. \emph{Acta Mathematica} 63, 193--248.

Li, D. (2019). Turbulent Prandtl number in the atmospheric boundary layer --- where are we now? \emph{Atmospheric Research} 216, 86--105. doi:10.1016/j.atmosres.2018.09.015

Monin, A. S., \& Obukhov, A. M. (1954). Basic laws of turbulent mixing in the surface layer of the atmosphere. \emph{Trudy Akademii Nauk SSSR Geofizicheskogo Instituta} 24, 163--187.

NEON (National Ecological Observatory Network) (2020). Bundled data products --- eddy covariance (DP4.00200.001), site WREF. https://data.neonscience.org

Ohana, R., et al. (2024). The Well: a large-scale collection of diverse physics simulations for machine learning. \emph{Advances in Neural Information Processing Systems (Datasets and Benchmarks Track)}.

Parish, E. J., \& Duraisamy, K. (2017). Non-Markovian closure models for large eddy simulations using the Mori--Zwanzig formalism. \emph{Physical Review Fluids} 2, 014604. doi:10.1103/PhysRevFluids.2.014604

Saifelden, A. (2026). What single-eddy-diffusivity closure discards: a Mori--Zwanzig audit, a scalar-independent memory time, and a divergence-cleaning design window. Companion manuscript, under editorial consideration at Physical Review (permanent submission ID assigned). \textbf{[P1]}

Smagorinsky, J. (1963). General circulation experiments with the primitive equations: I. The basic experiment. \emph{Monthly Weather Review} 91, 99--164.

Stull, R. B. (1988). \emph{An Introduction to Boundary Layer Meteorology}. Kluwer, Dordrecht.

\end{document}
